\section{TỔNG QUAN ĐỀ TÀI}
\subsection{Lý do chọn đề tài}
Hiện nay, có rất nhiều doanh nghiệp sở hữu mô hình gồm một trụ sở chính và nhiều chi nhánh phân bố ở các khu vực khác nhau. Điều này làm cho việc quản lý,cấu hình và giám sát hệ thống mạng của doanh nghiệp trở nên phức tạp hơn. Mô hình mạng truyền thống và cách trị cũ khiến doanh nghiệp gặp nhiều khó khăn trong việc dồng bộ cấu hình từng thiết bị, quản lý tập trung và giấm sát tình trạng hoạt động của hệ thống mạng.

Sự phát triển của công nghệ Cloud Computing và Software Definedd Networking (SDN) đã mở ra hướng mới trong việc quản lý hạ tầng mạng.Cloud Computing cung cấp môi trường triển khai linh hoạt, dễ mở rộng và tối ưu chi phí đầu tư, trong khi SDN tách biệt mặt phẳng điều khiển (Control Plane) khỏi mặt phẳng chuyển tiếp dữ liệu (Data Plane), cho phép quản lý tập trung toàn bộ thiết bị mạng thông qua một bộ điều khiển (SDN Controller). Nhờ đó, việc cấu hình VLAN, thiết lập Flow Rule, điều khiển lưu lượng và triển khai thiết bị mới có thể được thực hiện tự động, nhanh chóng và nhất quán.

Xuất phát từ những yêu cầu thực tế trên, nhóm chọn đề tài "Quản lý mạng doanh nghiệp trên nền tảng Cloud sử dụng SDN" nhằm nghiên cứu và xây dựng một mô hình quản lý mạng doanh nghiệp hiện đại,góp phần nâng cao hiệu quả quản trị mạng, giảm thời gian triển khai thiết bị, tăng khả năng mở rộng và đáp ứng xu hướng phát triển của các doanh nghiệp

\subsection{Mục tiêu thực hiện đề tài}
Đề tài hướng đến việc nghiên cứu, thiết kế và triển khai mô hình quản lý mạng doanh nghiệp trên nền tảng Cloud kết hợp với công nghệ Software Defined Networking (SDN), nhằm nâng cao hiệu quả quản trị và tự động hóa quá trình vận hành hệ thống mạng.

\subsection{Đối tượng và phạm vi nghiên cứu}
\subsection{Phương pháp nghiên cứu}
\subsection{Ý nghĩa thực tiễn của đề tài}
Đề tài được nghiên cứu và triển khai nhằm hỗ trợ doanh nghiệp giảm đáng kể thời gian cấu hình thiết bị, đơn giản hóa công tác quản trị mạng và tăng khả năng mở rộng khi bổ sung chi nhánh hoặc người dùng mới. Việc tích hợp giao diện Web giúp người quản trị dễ dàng theo dõi trạng thái hoạt động của toàn bộ hệ thống theo thời gian thực, đồng thời nhanh chóng phát hiện và xử lý các sự cố phát sinh.
\subsection{Bố cục báo cáo}
\begin{itemize}
     Nội dung báo cáo được trình bày trong năm chương như sau:

    \item Chương 1: Tổng quan đề tài. Trình bày lý do chọn đề tài, mục tiêu nghiên cứu, đối tượng và phạm vi nghiên cứu, phương pháp nghiên cứu, ý nghĩa thực tiễn và bố cục của báo cáo.
    \item Chương 2: Cơ sở lý thuyết. Giới thiệu các kiến thức nền tảng liên quan đến Cloud Computing, Private Cloud, Software Defined Networking (SDN), OpenFlow, SDN Controller, VLAN, VPN/GRE Tunnel và các công nghệ được sử dụng trong đề tài.
    \item Chương 3: Phân tích và thiết kế hệ thống. Phân tích yêu cầu hệ thống, thiết kế mô hình mạng doanh nghiệp, kiến trúc Cloud, SDN, cơ sở dữ liệu và giao diện Web quản lý.
    \item Chương 4: Triển khai hệ thống và đánh giá. Trình bày quá trình triển khai các thành phần của hệ thống, xây dựng Website quản lý, thực hiện các kịch bản kiểm thử và đánh giá hiệu năng dựa trên các tiêu chí đã đề ra.
\end{itemize}